\chapter{\IfLanguageName{dutch}{Stand van zaken}{State of the art}}%
\label{ch:stand-van-zaken}

% Tip: Begin elk hoofdstuk met een paragraaf inleiding die beschrijft hoe
% dit hoofdstuk past binnen het geheel van de bachelorproef. Geef in het
% bijzonder aan wat de link is met het vorige en volgende hoofdstuk.

% Pas na deze inleidende paragraaf komt de eerste sectiehoofding.

% Dit hoofdstuk bevat je literatuurstudie. De inhoud gaat verder op de inleiding, maar zal het onderwerp van de bachelorproef *diepgaand* uitspitten. De bedoeling is dat de lezer na lezing van dit hoofdstuk helemaal op de hoogte is van de huidige stand van zaken (state-of-the-art) in het onderzoeksdomein. Iemand die niet vertrouwd is met het onderwerp, weet nu voldoende om de rest van het verhaal te kunnen volgen, zonder dat die er nog andere informatie moet over opzoeken \autocite{Pollefliet2011}.

% Je verwijst bij elke bewering die je doet, vakterm die je introduceert, enz.\ naar je bronnen. In \LaTeX{} kan dat met het commando \texttt{$\backslash${textcite\{\}}} of \texttt{$\backslash${autocite\{\}}}. Als argument van het commando geef je de ``sleutel'' van een ``record'' in een bibliografische databank in het Bib\LaTeX{}-formaat (een tekstbestand). Als je expliciet naar de auteur verwijst in de zin (narratieve referentie), gebruik je \texttt{$\backslash${}textcite\{\}}. Soms is de auteursnaam niet expliciet een onderdeel van de zin, dan gebruik je \texttt{$\backslash${}autocite\{\}} (referentie tussen haakjes). Dit gebruik je bv.~bij een citaat, of om in het bijschrift van een overgenomen afbeelding, broncode, tabel, enz. te verwijzen naar de bron. In de volgende paragraaf een voorbeeld van elk.

% \textcite{Knuth1998} schreef een van de standaardwerken over sorteer- en zoekalgoritmen. Experten zijn het erover eens dat cloud computing een interessante opportuniteit vormen, zowel voor gebruikers als voor dienstverleners op vlak van informatietechnologie~\autocite{Creeger2009}.

% Let er ook op: het \texttt{cite}-commando voor de punt, dus binnen de zin. Je verwijst meteen naar een bron in de eerste zin die erop gebaseerd is, dus niet pas op het einde van een paragraaf.

Dit hoofdstuk bouwt verder op de probleemstelling uit de inleiding en biedt een diepgaand
overzicht van de huidige stand van zaken binnen het onderzoeksdomein. Achtereenvolgens
komen de connectiviteitsproblematiek in afgelegen gebieden, de fysische eigenschappen 
van radiofrequenties en atmosferische signaalpropagatie aan bod. Vervolgens worden 
multiplextechnieken, modulatieschema's en protocollen voor radiogebaseerde datacommunicatie 
geanalyseerd. Tot slot worden de technische trade-offs van straalverbindingen in zowel 
het gelicentieerde als het licentievrije spectrum (inclusief WiFi- en millimetergolfbanden) 
en de wettelijke randvoorwaarden voor radioamateurs in België besproken. Dit theoretisch kader 
vormt het fundament voor de methodologie en de eigen empirische realisatie in de 
volgende hoofdstukken.

\section{Connectiviteit in afgelegen gebieden en noodsituaties}
\label{sec:Connectiviteit in afgelegen gebieden en noodsituaties}

De beschikbaarheid van digitale bestanden is in de hedendaagse samenleving een kritieke
vereiste, maar blijft in uiteenlopende situaties een aanzienlijke uitdaging. In maritieme
omgevingen, afgelegen regio's en rampgebieden is betrouwbare internetconnectiviteit
structureel onvoldoende of volledig afwezig~\autocite{Xia2022}. \textcite{Zaballos2022}
beschrijven hoe noodsituaties zoals overstromingen, bosbranden en terroristische aanslagen
bijzonder hoge eisen stellen aan communicatiesystemen op het vlak van flexibiliteit,
mobiliteit, betrouwbaarheid en snelheid van inzet. Tegelijkertijd zijn net in dergelijke
situaties de conventionele communicatie-infrastructuren het meest kwetsbaar: mobiele
netwerken overbelasten, kabelinfrastructuur valt weg en satellietdiensten zijn niet altijd
beschikbaar of betaalbaar~\autocite{Zaballos2022}.

Commerciële satellietoplossingen bieden weliswaar globale dekking, maar kampen met hoge
abonnementskosten en aanzienlijke latency, waardoor ze voor particulieren en kleine
organisaties vaak financieel of praktisch ontoereikend zijn~\autocite{Xia2022}.
\textcite{Chaoub2021} argumenteren dat draadloze connectiviteit in afgelegen gebieden
een sleutelrol zal spelen in toekomstige communicatiegeneraties, waarbij de nadruk ligt
op robuuste, gedistribueerde systemen die zonder vaste infrastructuur kunnen
functioneren. Radiogebaseerde communicatie wordt in de literatuur herhaaldelijk naar
voren geschoven als een veelbelovend alternatief, mede omdat radioamateurs en
hulpverleningsorganisaties al decennialang gebruikmaken van radioverbindingen als
ruggengraat voor noodcommunicatie~\autocite{WinlinkEmcomm}. De International Telecommunication 
Union~\autocite{ITU2021} onderstreept in haar beleidslijnen dat het strategisch inzetten 
van digitale technologieën en radiocommunicatie de impact van rampen drastisch kan inperken 
en de continuïteit van de informatievoorziening waarborgt.

\section{Radiofrequenties en atmosferische signaalpropagatie}
\label{sec:Radiofrequenties en signaalpropagatie}

Om data succesvol via radiogolven te transporteren, is een fundamenteel begrip van de 
fysische laag van computercommunicatie essentieel, waarbij elektromagnetische golven worden 
gebruikt om bitstromen door de ether te sturen~\autocite{Tanenbaum2021}. De capaciteit 
van dit transportkanaal is onderhevig aan fundamentele theoretische limieten. \textcite{Hartley1928} 
toonde aan dat de maximale hoeveelheid informatie rechtstreeks evenredig is met de beschikbare 
bandbreedte en de transmissietijd. Dit concept werd geformaliseerd door \textcite{Shannon1948}, 
wiens wet met betrekking tot kanaalcapaciteit ($C$) aantoont dat de bovengrens van foutloze 
dataoverdracht afhankelijk is van de bandbreedte ($W$) en de signaal-ruisverhouding ($S/N$):

$$C = W \log_2 \left(1 + \frac{S}{N}\right)$$

In \textcite{Shannon1948PartII} wordt dieper ingegaan op de impact van continue kanalen 
met ruis, wat direct relevant is voor radiogebaseerde systemen. In de praktijk is 
signaalpropagatie sterk afhankelijk van de gekozen frequentieband en veranderlijke atmosferische 
omstandigheden~\autocite{Di Lorenzo2018}. Binnen de lagere frequentiebanden, zoals de kortegolf (HF), 
wordt het bereik bepaald door ionosferische omstandigheden die variëren naargelang het uur van de dag, 
het seizoen en de zonnecyclus. Een specifiek fenomeen binnen de hogere HF-banden (zoals 10 meter) 
en lagere VHF-banden (zoals 6 meter) is *Sporadic-E*. Hierbij ontstaan plotseling sterk 
geïoniseerde wolken in de E-laag van de ionosfeer, waardoor signalen met relatief lage vermogens 
over extreme afstanden gereflecteerd kunnen worden. Hoewel dit toestaat om duizenden kilometers 
te overbruggen, zijn deze condities inherent onvoorspelbaar en seizoensgebonden.

Voor kortere afstanden op hogere frequenties, zoals de VHF- en UHF-banden (bijvoorbeeld 70 centimeter), 
is de propagatie in principe beperkt tot de visuele horizon, tenzij er zich *troposferische propagatie* voordoet. Door specifieke meteorologische omstandigheden, zoals temperatuurinversies, ontstaan er 
brekingsindices in de lagere atmosfeer (*troposferische ducting*). Dit fenomeen stelt breedbandige 
radiosignalen in staat om de kromming van de aarde te volgen en honderden kilometers af te leggen. 
In tegenstelling tot ionosferische HF-reflecties vertoont troposferische propagatie in gematigde 
klimaten een relatief hoge en stabiele regelmaat, waardoor het bruikbaar is voor semi-permanente 
datacommunicatie over middellange afstanden.

\section{Multiplex- en modulatietechnieken in datacommunicatie}
\label{sec:Multiplextechnieken in radiogebaseerde datacommunicatie}

Wanneer meerdere knooppunten binnen een radionetwerk gelijktijdig gebruikmaken van het 
beschikbare spectrum, moeten multiplextechnieken worden toegepast om interferentie te voorkomen. 
\textcite{Tanenbaum2021} beschrijven hoe traditionele technieken zoals *Frequency Division Multiplexing* (FDM) en *Time Division Multiplexing* (TDM) het fysieke kanaal opdelen in aparte frequentiebanden of 
tijdsloten. Binnen moderne draadloze netwerkinfrastructuren wordt de capaciteit verder vergroot via 
spectrumaggregatie, waarbij niet-aaneengesloten frequentiebanden dynamisch worden gecombineerd 
tot één logisch kanaal~\autocite{MicrowavePtP5G}.

Bij het overdragen van data via atmosferische kanalen speelt de keuze van het modulatieschema een 
sleutelrol. Binnen toepassingen zoals Digitale Amateur Televisie (D-ATV) wordt vaak gewerkt met 
single-carrier modulatietechnieken zoals *Quadrature Phase Shift Keying* (QPSK) binnen de DVB-S2 
standaard om interfaces te simplificeren. Voor pure, robuuste IP-datacommunicatie over instabiele 
radiokanalen is *Orthogonal Frequency Division Multiplexing* (OFDM) echter superieur. OFDM verdeelt 
de datastroom over een groot aantal loodrechte subdraaggolven. Dit maakt het signaal extreem 
resistent tegen *multi-path fading* en looptijdverschillen (*delay spread*), fenomenen die inherent 
aanwezig zijn bij lange-afstandsverbindingen die gebruikmaken van atmosferische reflecties of 
troposferische breking.

\section{Protocollen voor radiogebaseerde dataoverdracht}
\label{sec:Protocollen voor radiogebaseerde dataoverdracht}

Waar de fysische laag de radiogolven beheert, bepalen de protocollen op de hogere lagen van het 
OSI-model hoe data logisch geordend en foutloos overgedragen wordt~\autocite{Tanenbaum2021}. 
Binnen de amateurradio- en noodcommunicatiewereld is *packet radio* historisch gezien de standaard 
voor data-over-radio~\autocite{ICTPPacketRadio}. Packet radio maakt gebruik van het AX.25-protocol 
om data op te delen in frames, wat efficiënt hergebruik van de frequentie mogelijk maakt en bovendien 
toestaat om het TCP/IP-protocol in te kapselen~\autocite{QSLPacketTCPIP}. Hierdoor kunnen standaard 
internetapplicaties draaien via het zogeheten AMPRNet (*Amateur Packet Radio Network*).

Voor grootschalige, gedistribueerde noodcommunicatie heeft het Winlink-ecosysteem zich ontwikkeld 
tot de wereldwijde standaard~\autocite{WinlinkOverview}. Winlink fungeert als een wereldwijd 
radio-e-mailsysteem dat met name zijn kracht ontleent aan het *Winlink Hybrid Network*~\autocite{WinlinkHybrid}. 
Dit netwerk combineert klassieke internetgateways met een autonoom radio-omleidingssysteem. 
Wanneer lokale internetinfrastructuur uitvalt, schakelt het systeem naadloos over op HF- of 
VHF/UHF-verbindingen om berichten via naburige, nog operationele radio-relayservers (RMS) te 
routeren~\autocite{WinlinkEmcomm}. Hoewel systemen zoals Winlink uitermate robuust zijn, zijn hun 
doorvoersnelheden door de beperkte bandbreedte laag. Voor applicaties die veel data verbruiken, 
zoals gecentraliseerde cloudomgevingen, zijn breedbandige alternatieven noodzakelijk~\autocite{Creeger2009}.

\section{Straalverbindingen: Privileges, ISM-banden en Millimetergolven}
\label{sec:Straalverbindingen als proof-of-concept}

Wanneer grote hoeveelheden data over middellange tot lange afstanden getransporteerd moeten worden 
zonder fysieke kabels, bieden point-to-point (PtP) straalverbindingen een onmisbaar, kosteneffectief 
en snel uit te rollen alternatief voor glasvezel~\autocite{AcomeMicrowaveFibre}. Moderne breedbandige 
straalverbindingen zijn in staat om gigabit-capaciteiten te overbruggen met minimale 
latency~\autocite{HighCapacityLongDistance}. Bij het ontwerp van zo'n netwerkarchitectuur moet 
een fundamentele afweging worden gemaakt tussen het gebruik van exclusieve radioamateurfrequenties 
en licentievrije spectrumalternatieven.

Het gebruik van de amateurbanden biedt unieke privileges, zoals hogere toegestane vermogens en 
specifieke propagatievoordelen, maar vereist dat elke gebruiker in het bezit is van een geldige 
radioamateurvergunning. Bovendien zijn deze banden in de praktijk gevoelig voor onvoorziene storingen 
van medegebruikers. Als alternatief kan worden uitgeweken naar licentievrije ISM-banden (*Industrial, 
Scientific and Medical*) op 2.4 GHz en 5.6 GHz via standaard WiFi-protocollen (zoals IEEE 802.11)~\autocite{LongRangePtP5GHz}. 
Hoewel deze hardware goedkoop en breed beschikbaar is, kampen deze banden met zware spectrumcongestie 
in dichtbevolkte gebieden.

Om aan deze congestie te ontsnappen en grotere bandbreedtes te realiseren, kan de licentievrije 
60 GHz-band (de V-band) worden ingezet, waar zeer brede kanalen beschikbaar zijn voor gigabit-doorvoer. 
De bruikbaarheid hiervan voor lange afstanden wordt echter fundamenteel begrensd door atmosferische 
fysica. Rond de 60 GHz treedt een sterke elektromagnetische resonantie op met zuurstofmoleculen ($O_2$), 
wat resulteert in een specifieke piek in atmosferische absorptie (tot wel 15 dB/km). Zelfs bij een 
perfecte optische zichtverbinding (*Line-of-Sight*) beperkt deze moleculaire demping het bereik 
van 60 GHz-straalverbindingen inherent tot zeer korte afstanden, waardoor ze ongeschikt zijn voor 
interstedelijke backhaullinks~\autocite{Di Lorenzo2018}.

\section{Wettelijk kader en regelgeving in België}
\label{sec:Wettelijk kader en regelgeving in België}

Elk technisch ontwerp dat gebruikmaakt van het radiospectrum moet strikt voldoen aan de nationale 
en internationale wetgeving om schadelijke interferentie met primaire diensten te voorkomen. In 
België wordt het gebruik van radiofrequenties door radioamateurs formeel gereguleerd door het 
Belgisch Instituut voor Postdiensten en Telecommunicatie~\autocite{BIPT2019}. Het wettelijk besluit 
definieert nauwkeurig welke frequentiebanden openstaan voor de verschillende amateurlicenties, 
welke transmissiemodi zijn toegestaan, en begrenst strikt het maximaal toegestane zendvermogen.

Voor de praktische vertaling van deze wetgeving naar werkbare bandplannen vertrouwen Belgische 
radioamateurs op de richtlijnen van de Unie van Belgische Zendamateurs~\autocite{UBAFrequenties}. 
De UBA stemt de nationale frequentieverdelingen af met de International Amateur Radio Union (IARU) 
en de overheid om ervoor te zorgen dat experimentele datacommunicatie binnen de wettelijk 
vastgelegde grenzen blijft. Een succesvolle proof-of-concept moet derhalve expliciet ontworpen worden 
binnen de vermogens- en frequentiebeperkingen die door deze instanties zijn opgelegd~\autocite{BIPT2019, UBAFrequenties}.

% \begin{figure}
%   \centering
%   \includegraphics[width=0.8\textwidth]{grail.jpg}
%   \caption[Voorbeeld figuur.]{\label{fig:grail}Voorbeeld van invoegen van een figuur. Zorg altijd voor een uitgebreid bijschrift dat de figuur volledig beschrijft zonder in de tekst te moeten gaan zoeken. Vergeet ook je bronvermelding niet!}
% \end{figure}

% \begin{listing}
%   \begin{minted}{python}
%     import pandas as pd
%     import seaborn as sns

%     penguins = sns.load_dataset('penguins')
%     sns.relplot(data=penguins, x="flipper_length_mm", y="bill_length_mm", hue="species")
%   \end{minted}
%   \caption[Voorbeeld codefragment]{Voorbeeld van het invoegen van een codefragment.}
% \end{listing}

% \lipsum[7-20]

% \begin{table}
%   \centering
%   \begin{tabular}{lcr}
%     \toprule
%     \textbf{Kolom 1} & \textbf{Kolom 2} & \textbf{Kolom 3} \\
%     $\alpha$         & $\beta$          & $\gamma$         \\
%     \midrule
%     A                & 10.230           & a                \\
%     B                & 45.678           & b                \\
%     C                & 99.987           & c                \\
%     \bottomrule
%   \end{tabular}
%   \caption[Voorbeeld tabel]{\label{tab:example}Voorbeeld van een tabel.}
% \end{table}

